\chapter{Container, Kubernetes, dan Serverless}

\section*{Tujuan Pembelajaran}
\begin{itemize}
  \item Menjelaskan perbedaan \textit{virtual machine}, \textit{container}, dan \textit{serverless} dengan analogi dan contoh.
  \item Menjelaskan cara Kubernetes menjaga sistem tetap sesuai keadaan yang diinginkan.
  \item Memilih \textit{runtime} untuk setiap kandidat layanan berdasarkan kriteria yang eksplisit.
\end{itemize}

\section{Pendahuluan}

% TODO
% Cakupan: aplikasi yang berjalan baik di satu rumah sakit sering gagal di rumah sakit lain karena lingkungannya berbeda; pada saat yang sama, sebagian layanan hanya dipakai beberapa menit sehari tetapi servernya menyala terus.

\section{Virtual Machine dan Container}

% TODO
% Cakupan: VM dan container dengan analogi peti kemas; container image dan registry; manfaat konsistensi lingkungan; batas isolasi container dibandingkan VM.
% Pendalaman teknis: menulis Dockerfile untuk aplikasi contoh dan menyimpan image di registry.

\section{Orkestrasi Container dan Konsep Kubernetes}

% TODO
% Cakupan: mengapa container perlu orkestrasi; konsep desired state dan rekonsiliasi; pod, deployment, service, self-healing, dan autoscaling pada tingkat konsep; data survei CNCF 2025 tentang adopsi Kubernetes di produksi dan sebagai platform workload AI.
% Pendalaman teknis: menjalankan aplikasi contoh di klaster Kubernetes lokal dengan kind atau minikube.

\section{Managed Kubernetes dan PaaS}

% TODO
% Cakupan: tanggung jawab yang diambil alih penyedia pada managed Kubernetes; PaaS sebagai alternatif yang lebih sederhana; biaya keterampilan dan operasi yang sering diremehkan.
% Kasus Medikara: apakah tim TI Medikara Health yang kecil sanggup mengelola Kubernetes sendiri.

\section{Serverless dan Function as a Service}

% TODO
% Cakupan: serverless dan FaaS; skala sampai nol dan biaya per eksekusi; cold start, batas waktu eksekusi, dan ketergantungan pada layanan penyedia; pola beban yang cocok.
% Kasus Medikara: pengingat jadwal kontrol pasien dan pemrosesan hasil laboratorium yang datang berkala.

\section{Memilih Runtime dan Mengelola Vendor Lock-in}

% TODO
% Cakupan: kriteria memilih runtime (pola beban, latensi, keterampilan tim, biaya, lokasi data); matriks keputusan dari SaaS sampai VM; vendor lock-in, portabilitas, dan biaya keluar dari penyedia; kaitan dengan matriks lokasi data Bab 2.

\section{Aktivitas: Matriks Keputusan Runtime}

% TODO
% Tujuan aktivitas: menetapkan runtime untuk setiap kandidat layanan dengan kriteria yang dituliskan lebih dulu.
% Peran dalam tim: anggota teknis menjelaskan konsekuensi operasional setiap opsi; anggota non-teknis menilai biaya, ketersediaan tenaga kerja di pasar, dan risiko ketergantungan pada penyedia; seluruh anggota menyepakati bobot kriteria.
% Langkah: (1) ambil kandidat layanan dari Bab 6 dan 7; (2) tetapkan kriteria dan bobotnya; (3) nilai opsi SaaS, PaaS, managed Kubernetes, serverless, dan VM; (4) pilih runtime dan catat alasannya dalam ADR singkat.
% Luaran dan tonggak capstone: C3, matriks keputusan runtime per layanan.
% Berkas sources yang direncanakan: sources/chapter08/matriks-keputusan-runtime.xlsx, sources/chapter08/contoh-container/ (opsional, Dockerfile dan manifest Kubernetes)

\section{Ringkasan}

% TODO

\section{Latihan}

\begin{enumerate}
  \item Jelaskan kepada direktur keuangan mengapa layanan pengingat jadwal kontrol pasien cocok dijalankan secara \textit{serverless}.
  \item Kapan Kubernetes justru menjadi beban bagi organisasi? Berikan dua kondisi beserta alasannya.
  \item Bandingkan risiko \textit{vendor lock-in} pada layanan \textit{serverless} dan pada \textit{managed} Kubernetes.
\end{enumerate}
